A set of data containing the apparent positions of 4 Solar System objects
over a period of 1 calendar year is given in Table 1. Show your method of
data analysis carefully and answer the following questions.

\begin{center}
\begin{tabular}{lll}
Location of observer & Latitude: & N $18^\circ 47' 00.0''$ \\
 & Longitude: & E $98^\circ 59' 00.0''$ \\
\end{tabular}
\end{center}

% FIGURE NEEDED: Table 1 — apparent positions of the 4 objects (A, B, C, D)
% over one calendar year. Not present in the source paper (2007_da.pdf);
% it was distributed separately with the exam materials.

\begin{description}
\item[3.1] Put the letters A, B, C and D beside the appropriate objects on
  the answer sheet. \textbf{(2 points)}
\item[3.2] During the period of observation, which object could be observed
  for the longest duration at night time? \textbf{(1 point)}
\item[3.3] What was the date corresponding to the situation in 3.2?
  \textbf{(1 point)}
\item[3.4] Assuming the orbits are coplanar (lie on the same plane) and
  circular, indicate the positions of the four objects and the Earth on the
  date in 3.3, in the orbit diagram provided in your answer sheet. The answer
  (sheet) must show one of the objects as the Sun at the centre of the Solar
  System. Other objects including the Earth must be specified together with
  the correct values of elongation on that date. \textbf{(4 points)}
\end{description}
